% GENERATED FILE. Edit canonical lessons, then run npm run generate:books.
% canonical-source-hash: fnv1a64:f55741a1fc4a833c
% canonical-lessons: MW-C151-ar-together, MW-C151-ar-choice, MW-C151-pan-contrast, MW-C151-kyunkai-reason

\chapter{And, Or, But, and One Reason}
\label{ch:mw-a1-joining-choice}
\hlchaptermodality{151}

\begin{chapteropening}
\textbf{By the end of this chapter:} \emph{I can join two familiar things, offer a choice, contrast two short statements, and give one reason.}

Join a pair, offer tea or water, contrast hot tea with cold water, and give one short reason with kyunkai.
\end{chapteropening}

\section[\texorpdfstring{\mw{अर} (ar)}{अर (ar)}]{\mw{अर} (ar) — put two familiar things together}

\label{lesson:MW-C151-ar-together}



\begin{quote}
Say \textbf{\mw{कुत्तो}} (\emph{kutto}, a dog). Then say \textbf{\mw{बिल्ली}} (\emph{billī}, a cat). Keep the two words separate for one more breath.
\end{quote}

\subsection*{You'll want to know: \mw{अर}}
\begin{quote}
\textbf{\mw{अर}} — \emph{ar} — \textbf{and}, joining this thing with that thing
\end{quote}

The Government of India's \emph{Linguistic Survey of India — Rajasthan, Part I} records Marwari \textbf{\mw{अर}} as the particle that links words or groups of words. Its own example joins a dog and a cat. This lesson keeps the book's already-taught word \textbf{\mw{बिल्ली}} for cat.

\begin{cousinweb}
The official survey attests the Marwari form and its “and / or” use, but does not derive it. Its origin remains unknown in the source used for this lesson.
\end{cousinweb}

\begin{grammarlens}[title={What you've built}]
Put \textbf{\mw{अर}} between the two things:

\begin{quote}
\textbf{\mw{कुत्तो} \mw{अर} \mw{बिल्ली}} — \emph{kutto ar billī} — a dog and a cat
\end{quote}

Nothing changes on either noun. The small joining word does the new work.
\end{grammarlens}

\subsection*{Your turn}
\begin{itemize}
\raggedright
  \item \emph{Say it:} \textbf{\mw{कुत्तो} \mw{अर} \mw{बिल्ली}}
  \item \emph{Switch:} cat first — \textbf{\mw{बिल्ली} \mw{अर} \mw{कुत्तो}}
  \item \emph{Recall:} say only the joining word — \textbf{\mw{अर}}
\end{itemize}

\begin{tcolorbox}[breakable,colback=teal!4,colframe=teal!35!black,title={Before you move on}]
Say \textbf{\mw{कुत्तो} \mw{अर} \mw{बिल्ली}} once without looking. What is the one new piece? (\textbf{\mw{अर}}, \emph{ar}, and.)

Source: Government of India, \emph{Linguistic Survey of India — Rajasthan, Part I}, section 4.5, “Conjunctive particle”, pp. 125–126.
\end{tcolorbox}
\section[\texorpdfstring{\mw{चाय} \mw{अर} \mw{पाणी}? (chāy ar pāṇī?)}{चाय अर पाणी? (chāy ar pāṇī?)}]{\mw{चाय} \mw{अर} \mw{पाणी}? — offer two familiar choices}

\label{lesson:MW-C151-ar-choice}



\begin{quote}
Say \textbf{\mw{चाय}} (\emph{chāy}, tea), then \textbf{\mw{पाणी}} (\emph{pāṇī}, water). Join them once with \textbf{\mw{अर}} as “tea and water”.
\end{quote}

\begin{grammarlens}[title={What you've built}]
The official survey gives \textbf{\mw{अर}} both “and” and “or” value. When somebody offers two alternatives with a question contour, the same small word carries the choice:

\begin{quote}
\textbf{\mw{चाय} \mw{अर} \mw{पाणी}?} — \emph{chāy ar pāṇī?} — Tea or water?
\end{quote}

The setting and the rising question do the separating. Do not import a Hindi connector to make the choice look more familiar.
\end{grammarlens}

\begin{cousinweb}
This lesson adds a new use, not a new connector. The official survey attests \textbf{\mw{अर}} with “and / or” value but does not derive it; its deeper origin remains unknown in that source.
\end{cousinweb}

\subsection*{Your turn}
\begin{itemize}
\raggedright
  \item \emph{Ask:} \textbf{\mw{चाय} \mw{अर} \mw{पाणी}?}
  \item \emph{Answer:} \textbf{\mw{पाणी।}}
  \item \emph{Ask:} reverse the alternatives — \textbf{\mw{पाणी} \mw{अर} \mw{चाय}?}
  \item \emph{Answer:} \textbf{\mw{चाय।}}
\end{itemize}

One short item is enough for the answer. This lesson teaches how to hand over a real choice; it does not yet teach “I prefer”.

\begin{tcolorbox}[breakable,colback=teal!4,colframe=teal!35!black,title={Before you move on}]
Offer tea or water. Answer with one item. This closes a tiny two-turn choice, not a general preference statement.

Source: Government of India, \emph{Linguistic Survey of India — Rajasthan, Part I}, section 4.5, “Conjunctive particle”, pp. 125–126.
\end{tcolorbox}
\section[\texorpdfstring{\mw{पण} (paṇ)}{पण (paṇ)}]{\mw{पण} (paṇ) — set one statement against another}

\label{lesson:MW-C151-pan-contrast}



\begin{quote}
Say \textbf{\mw{चाय} \mw{गरम} \mw{है।}} Tea is hot. Then say \textbf{\mw{पाणी} \mw{ठंडो} \mw{है।}} Water is cold. Keep both short statements intact.
\end{quote}

\subsection*{You'll want to know: \mw{पण}}
\begin{quote}
\textbf{\mw{पण}} — \emph{paṇ} — \textbf{but, yet}
\end{quote}

The Government of India's survey records Marwari \textbf{\mw{पण}} as the particle for “yet/but”. A Marwari dictionary independently prints the same Devanagari form and labels it a conjunction meaning “but”.

\begin{cousinweb}
The two sources attest the Marwari form and meaning but do not give an etymology. Its origin remains unknown in the sources used for this lesson.
\end{cousinweb}

\begin{grammarlens}[title={What you've built}]
Keep the two familiar statements and put \textbf{\mw{पण}} at their turning point:

\begin{quote}
\textbf{\mw{चाय} \mw{गरम} \mw{है}, \mw{पण} \mw{पाणी} \mw{ठंडो} \mw{है।}}
\end{quote}

\begin{quote}
\emph{chāy garam hai, paṇ pāṇī ṭhaṇḍo hai.}
\end{quote}

\begin{quote}
Tea is hot, but water is cold.
\end{quote}

The first statement stays true. \textbf{\mw{पण}} tells the listener that the second one pulls in a different direction.
\end{grammarlens}

\subsection*{Your turn}
\begin{itemize}
\raggedright
  \item \emph{Say it:} \textbf{\mw{चाय} \mw{गरम} \mw{है}}
  \item \emph{Add:} \textbf{\mw{पण} \mw{पाणी} \mw{ठंडो} \mw{है}}
  \item \emph{Say it:} the whole contrast once, keeping \textbf{\mw{पण}} in the middle
\end{itemize}

\begin{tcolorbox}[breakable,colback=teal!4,colframe=teal!35!black,title={Before you move on}]
What single word marks the turn? (\textbf{\mw{पण}}, \emph{paṇ}, but/yet.) Say the whole contrast once more.

Sources: Government of India, \emph{Linguistic Survey of India — Rajasthan, Part I}, section 4.5, pp. 125–126; \emph{Marwari–Urdu–English Dictionary}, entry \textbf{\mw{पण}}.
\end{tcolorbox}
\section[\texorpdfstring{\mw{बाटी} \mw{गरम} \mw{है} \mw{क्यूंकै} … (bāṭī garam hai …)}{बाटी गरम है क्यूंकै … (bāṭī garam hai …)}]{\mw{क्यूंकै} — give one short reason}

\label{lesson:MW-C151-kyunkai-reason}



\begin{quote}
Bring back \textbf{\mw{क्यूंकै}} (\emph{kyū̃kai}, because), \textbf{\mw{बाटी}} (\emph{bāṭī}), \textbf{\mw{तवो}} (\emph{tavo}, a griddle), and \textbf{\mw{गरम} \mw{है}} (\emph{garam hai}, is hot).
\end{quote}

\begin{cousinweb}
Chapter 117 takes \textbf{\mw{क्यूंकै}} apart as \emph{kyū̃} plus \emph{kai}. CDIAL lists Old Marwari \emph{kyu} under Sanskrit \emph{kim} “what” (CDIAL 3164); the origin of the second part remains not settled. This lesson reuses that word and adds only the two-clause reason frame.
\end{cousinweb}

\begin{grammarlens}[title={What you've built}]
Chapter 117 taught the word \textbf{\mw{क्यूंकै}}. A word meaning “because” is not yet a reason. Put one complete familiar statement before it and another after it:

\begin{quote}
\textbf{\mw{बाटी} \mw{गरम} \mw{है} \mw{क्यूंकै} \mw{तवो} \mw{गरम} \mw{है।}}
\end{quote}

\begin{quote}
\emph{bāṭī garam hai kyū̃kai tavo garam hai.}
\end{quote}

\begin{quote}
The bati is hot because the griddle is hot.
\end{quote}

The first part says what is so. The second gives one reason. Both small statements keep \textbf{\mw{है}} at their ends.
\end{grammarlens}

\subsection*{Your turn}
\begin{itemize}
\raggedright
  \item \emph{Say it:} \textbf{\mw{बाटी} \mw{गरम} \mw{है}}
  \item \emph{Add:} \textbf{\mw{क्यूंकै} \mw{तवो} \mw{गरम} \mw{है}}
  \item \emph{Say it:} the complete reason line once
\end{itemize}

Now retrieve the chapter in four small moves:

\begin{enumerate}
\raggedright
  \item join \textbf{\mw{कुत्तो}} and \textbf{\mw{बिल्ली}} with \textbf{\mw{अर}};
  \item offer \textbf{\mw{चाय} \mw{अर} \mw{पाणी}?} and answer with one;
  \item contrast hot tea with cold water using \textbf{\mw{पण}}; and
  \item give the bati-and-griddle reason using \textbf{\mw{क्यूंकै}}.
\end{enumerate}

\begin{tcolorbox}[breakable,colback=teal!4,colframe=teal!35!black,title={Before you move on}]
Say the whole reason once without looking. This is one small curriculum frame against a partial inventory: it is not an A1 mock, not A1 readiness, not an attainment claim, and not human validation.

Source for \textbf{\mw{क्यूंकै}}: the existing chapter-117 Marwadi lesson and its cited dictionary/CDIAL trail. No new lexical source is claimed here.
\end{tcolorbox}
